\chapter{绪论}

\section{研究背景与意义}

\subsection{大学生心理健康问题的现状}

进入高等教育阶段之后，学生实际要面对的压力来源，往往比外界所设想的更加复杂。学业上的竞争、就业方面的压力、人际关系的调适、经济负担，以及对自我价值持续不断的追问，一同构成了当代大学生心理负荷的主要来源。近些年来，国内多项大规模调查都显示，高校学生心理健康问题的检出率正处在上升过程当中。按照中国科学院心理研究所发布的《中国国民心理健康发展报告（2021—2022）》来看，大学生群体中存在抑郁风险的比例约为 21.4\%，焦虑问题的比例同样不容乐观\cite{cass2023}。世界卫生组织的相关报告也指出，全球约有 20\% 的青少年存在心理健康问题，其中相当大的一部分人由于未能及时获得专业支持，进一步错过了干预时机\cite{who2022}。

心理健康方面的问题，往往会对学业表现、人际交往以及身体健康带来多方面且相对深远的影响。部分学生如果长期处于焦虑或者抑郁的状态当中，学习方面的效率往往会跟着下降，情况严重的时候，甚至还会出现退学或者休学之类的问题；更需要重视的是，要是心理困扰不断累积并且一直没有得到及时的疏导，那么还有可能进一步引发极端事件。也正是鉴于这些现实情况，高校开展心理健康工作的定位，已经由一种辅助性服务，逐步转向教育管理当中的核心议题之一。

\subsection{传统线下咨询的痛点}

现阶段，国内大多数高校都已经设立了心理咨询中心，并且配备了专职或者兼职咨询师，来为在校学生提供面对面的咨询服务。不过，这种传统的线下模式在实际开展运行的过程当中，仍旧显现出了一些比较明显的结构性不足。

\textbf{预约方面的流程相对比较繁琐。}学生通常需要到咨询中心现场办理，或者借助电话、邮件等方式来进行登记预约，排队等候的时间往往比较长，短的话需要数天，长的话则可能达到数周。这样的时间延后，对于正处在急性心理困扰阶段的学生而言，往往就意味着会错过最佳的干预窗口。

\textbf{寻求帮助的门槛相对还是偏高。}心理咨询在社会文化层面上，仍旧存在一定程度的污名化现象。不少学生会因为担心被同学或者老师知晓而产生顾虑，进而选择独自承受，而不是主动去寻求专业帮助。线下面谈这种模式本身并不具备匿名性，而这一心理壁垒往往会在一定程度上抑制学生主动求助的意愿\cite{li2017}。

\textbf{档案管理工作相对比较分散。}咨询记录、预约信息以及评估结果等数据，往往会以纸质材料或者相互分散的电子文件形式来进行存储，这样一来，不仅不利于咨询师对个案进展开展持续追踪，还会使学校管理层难以形成系统化的数据视图，进而对服务质量的持续改进形成制约。

\textbf{服务开放时间方面存在一定限制。}心理咨询中心一般只会在工作日的固定时段对外开展服务，因此较难覆盖学生在深夜、节假日等非工作时间所产生的情绪支持需求。

\subsection{构建线上平台的必要性与社会意义}

信息技术的持续发展，为化解前面这些困境提供了较为切实可行的实现路径。构建一套面向高校的心理健康咨询 Web 平台，不仅是技术层面的工程实践，同时也是对现实需求加以回应的一项社会价值行动。

从可及性的角度来讲，线上平台能够较为有效地把时间以及空间方面的限制打破，学生可以在任何地点、任何时段去发起预约，或者和 AI 助手开展对话，这样就会让求助的第一步更容易被迈出去。从私密性的角度来讲，数字化渠道天然会形成一层心理上的距离感，有助于对由污名化所带来的顾虑进行减轻。从管理效能这一角度来看，采用结构化的数据存储方式之后，咨询师会更方便地去查阅既往记录，并开展案例笔记的撰写工作；学校管理者也能够凭借数据统计来了解整体服务态势，进而为资源配置以及政策决策提供相应依据。

需要说明的是，随着大语言模型技术逐步走向成熟，平台在引入 7×24 小时 AI 心理助手方面，已经具备了较为现实的条件。AI 助手并不是对专业咨询师的替代，而是把它当作情绪疏导以及初步支持的第一层触达；对于症状相对较轻、暂时还不需要人工干预的学生来说，这一功能可以在一定程度上对服务空白进行填补。

综合前文内容来看，构建一套把在线预约、AI 辅助对话以及数据管理整合于一体的心理健康咨询 Web 平台，对于提高高校心理服务整体效能、推动心理健康工作朝着数字化转型方向发展，都具有较为明确的理论价值和现实实践意义。


\section{国内外研究现状}

\subsection{国外心理健康平台的发展}

国外在线心理健康服务的起步时间相对比较早，并且在后续发展过程当中逐步形成了若干具有代表性的商业化平台，它们在实际运行方面所积累的经验，对本文的研究工作具有较强的参考价值。

\textbf{BetterHelp} 是目前全球范围内规模最大的在线心理咨询平台之一，于 2013 年在美国成立。该平台选用订阅制的运行模式，为用户去匹配持证心理咨询师，并且提供文字、语音以及视频这三种沟通方式，同时还允许用户按照自身偏好来自主切换或者更换咨询师。它的核心优势主要体现在对咨询师资源进行规模化整合以及沟通形式较为灵活这两个方面，不过对于一部分经济压力相对较大的用户而言，订阅费用依然会形成一定障碍\cite{baumel2019}。

\textbf{Talkspace} 也是美国发展得较为成熟的在线心理咨询平台之一，除了面向个人用户提供 B2C 服务之外，还把应用进一步拓展到 B2B 场景，并且和企业、学校以及保险机构等开展合作，提供员工援助计划（EAP）以及学生心理健康支持项目\cite{mohr2018}。这种借助机构合作来运行的模式，在高校心理服务场景当中有较强的直接借鉴意义。

\textbf{Woebot} 则体现了另一类技术路径，也就是把认知行为疗法（CBT）当作理论基础，借助对话式 AI 提供心理支持，并且不需要人工咨询师参与其中\cite{fitzpatrick2017}。Woebot 的相关实践说明，结构化 AI 对话在对轻度焦虑以及抑郁症状进行缓解方面能够起到积极作用，不过它的局限性也比较明显：如果面对中重度心理问题，AI 仍旧较难替代专业人工干预。

从整体情况来进行分析，国外平台在商业模式的创新、用户体验的设计以及 AI 辅助应用等方面仍走在前面，但它的服务体系大多还是面向成人个体用户，针对高校场景所开展的定制化建设依然比较不足，同时在数据合规要求以及本土化适配方面，和国内实际环境之间仍有较为明显的差距。

\subsection{国内高校心理咨询信息化进展}

国内心理健康信息化建设在系统化推进方面起步相对偏晚，但近些年在政策持续推动之下，已经呈现出较为明显的提速发展趋势。教育部于 2018 年印发《高等学校学生心理健康教育指导纲要》，明确提出高校要加强心理健康信息化建设，并且进一步把"测、评、咨、防"一体化服务体系加以完善\cite{moe2018}。在这一政策背景之下，一些高校开始选用商业化心理健康管理平台，或者委托软件公司进行定制化系统的开发工作。

目前国内高校心理咨询信息化系统在多数情况下仍把心理测评以及档案管理当作核心内容，通常有在线心理量表测试功能，比如 SCL-90、SDS、SAS 等，同时还包括预约登记、报告生成等基础功能\cite{chen2020}。这类系统虽然在数据积累以及管理规范化方面发挥了较为重要的作用，但也同时暴露出一些较为普遍的局限：

\textbf{AI 集成水平依旧相对较低。}绝大多数现有系统还没有把大语言模型或对话式 AI 真正纳入到服务流程当中，这就导致学生在非工作时段往往难以及时获得回应，系统因此更多只是充当"信息记录工具"，而并非真正意义上的"服务提供者"。

\textbf{用户体验方面整体上仍有不足。}部分系统在界面设计较为陈旧、操作流程偏长、前端技术栈老化以及移动端适配不够完善等方面存在问题，从而在一定程度上削弱了学生的使用意愿。

\textbf{各项功能之间仍存在相互分离的问题。}预约、咨询记录、测评数据以及公告通知通常分散在不同的子系统当中，缺少统一的数据视图，这就使得咨询师以及管理者在实际使用过程当中需要频繁地切换界面。

\subsection{现有研究的不足与本文的切入点}

从整体情况来看，当前国内外有关于这一领域的研究以及产品，在以下几个方面仍有进一步完善的空间：其一，专门面向高校场景、同时服务学生、咨询师以及管理者的三端协同平台依旧比较少见；其二，把实时流式 AI 对话以及传统预约管理进行深度整合的系统，目前仍然不算多见；其三，在面向信息技术专业的系统实现层面，能够供参考的开源资料相对较少，因此较难为同类系统后续的建设工作提供工程层面的借鉴。

也正是在上述背景之下，本文独立开展了一套面向高校场景的心理健康咨询 Web 平台的设计与实现工作，重点对三端统一架构、SSE 流式对话技术以及标准化 RESTful API 在这一场景中的工程化应用进行了探索。


\section{研究内容与组织结构}

\subsection{研究内容概述}

本文所开展的核心工作，是对一套较为完整的大学生心理健康咨询 Web 平台进行设计与实现，相关的研究内容主要包括以下几个方面：

\textbf{三端功能架构的设计与实现。}系统把学生端、咨询师端以及学校管理端划分成三个功能域，各端都围绕对应角色的核心业务需求来进行功能设计，并且借助统一的后端 API 以及角色权限中间件，实现数据共享以及权限隔离。

\textbf{AI 心理助手的集成方案。}平台接入 DeepSeek 大语言模型\cite{deepseek2024}，并借助 Server-Sent Events 技术来实现流式对话，使 AI 回复能够以逐词输出的形式进行呈现，从而对交互过程的流畅性进行提升。系统提示词经过有针对性的设计，使 AI 助手能够以专业、温和的方式回应学生的情绪倾诉，并自动维护最近 20 条消息的上下文记忆，以保证对话在前后衔接上的连贯性。

\textbf{安全认证以及权限控制机制。}系统借助 JSON Web Token 来实现无状态认证\cite{jones2015}，同时配合角色守卫中间件，对三端接口进行细粒度保护；用户密码会在经过 bcrypt 哈希加密之后再存储\cite{provos1999}，从而在数据层面降低明文泄露的风险。

\textbf{数据库设计以及编码适配。}系统数据库统一选用 utf8mb4 编码，并且围绕连接配置、Sequelize 模型定义以及表结构这三个层面开展一致化处理，从而确保中文内容得以实现无损存取。

\subsection{论文章节安排}

本文共分七章，各章内容安排如下：

第 1 章作为绪论部分，主要对大学生心理健康问题的现实状况以及研究背景进行阐述，同时对国内外相关研究以及产品进展加以回顾，并进一步明确本文的研究内容与技术路线。

第 2 章为相关技术与理念介绍，对系统所采用的主要技术栈进行简要梳理，涵盖 Vue 3、Node.js + Express、Sequelize、MySQL、JWT 以及 SSE 流式通信等方面，为后续章节提供技术背景。

第 3 章为系统需求分析部分，主要从功能需求、非功能需求以及角色权限这三个维度，对平台的整体需求进行系统性分析，进而形成需求规格说明。

第 4 章属于系统总体设计部分，主要包括系统架构设计、数据库设计以及功能模块划分等内容，为后续系统的实现工作提供整体蓝图。

第 5 章属于系统功能实现部分，介绍工程目录结构，并围绕用户认证、学生端、咨询师端、学校管理端以及 AI 心理助手等模块，结合关键代码，对核心实现细节进行阐述。

第 6 章为系统测试部分，围绕功能、性能以及安全性这三个方面开展相关测试，并对测试结果进行分析。

第 7 章属于总结与展望部分，主要对系统的整体情况加以归纳和评价，指出当前仍然存在的不足，同时对未来的改进方向进行展望。


\section*{本章小结}

本章结合大学生心理健康问题的现实情境，对传统线下咨询在预约效率、求助门槛以及档案管理等方面所存在的结构性不足进行了较为系统的分析；借助对国内外相关研究以及平台进展的梳理工作，明确了本研究把"三端协同架构 + SSE 流式 AI 对话"当作核心切入点；在这一基础之上，进一步说明了本研究的四项主要内容以及全文七章的组织安排。后续各章将在本章所奠定的研究基础之上，逐步开展系统的设计、实现以及验证工作。
